\section{Zadanie 12}
Czy każda funkcja ciągła $f: \mathbb{R} \to \mathbb{R}$ jest mierzalna?

\subsection*{Rozwiązanie}
Tak, pokażemy, że każda taka funkcja jest $\mathcal{B}(\mathbb{R})$-mierzalna.
Niech $\mathcal{A}_0$ oznacza rodzinę wszystkich zbiorów otwartych w $\mathbb{R}$. Wiemy, że $\sigma(\mathcal{A}_0) = \mathcal{B}(\mathbb{R})$
Na podstawie definicji Cauchy'ego ciągłości funkcji $f$ mamy
\[
\forall A \in \mathcal{A}_0 \quad f^{-1}(A) \in \mathcal{A}_0,
\]
czyli przeciwobraz każdego zbioru otwartego jest również zbiorem otwartym.

Teraz wykorzystamy zadanie 10 z zestawu 3. Podstawiając tam:
\[
\Omega _1 = \Omega _2 = \mathbb{R};  \Sigma _1 = \Sigma _2 = \mathcal{B}(\mathbb{R}) ; \mathcal{G} = \mathcal{A}_0
\]
otrzymujemy, że f jest mierzalna.

\subsection*{Odpowiedź}
\[
\boxed{\text{Tak, w ujęciu borelowskim}}
\]
